\documentclass[../main.tex]{subfiles}

\begin{document}

\section{Linear Algebra}

Linear algebra is a foundational building block of ML as majority
of the components are written in very high dimensional arrays
and tensors. Without exaggeration, training and using a model
can be considered as performing billions of linear algebra
operations. Therefore, it is encouraged to get familiar with
the notations as they will be used throughout the future notes.

\subsection{Foundations}

First, let's get started with definitions. What is a vector?
You probability have learned vectors as a quantity with both
magnitude and direction. The same definition applies in math
and ML, but how we represent and perceive them differs.

\begin{definition}{}{}
A \textbf{vector} is quantity with magnitude and directions.
The set of vectors that satisfies addition and multiplication
properties is called \textbf{vector space}. In math and ML,
vectors are often represented as coordinates in vector space.
\textbf{Dimensions} are the number of components that composes
a vector.
\end{definition}

Vectors with specific dimensions have dedicated names.

\begin{definition}{}{}
\textbf{Scalars} are 0-dimensional quantities. A 1-dimensional
vector is simply called \textbf{vector}. A 2D vector is called
\textbf{matrix}. For vectors with 3+ dimensions, the term
\textbf{tensor} is used.
\end{definition}

Now scalars are not considered vectors because it is only
defined with magnitude, without direction. Geometrically speaking,
the dimensions here can be interpreted as the dimension that
we usually say in science. For example, a matrix can be plotted
in a 2D space with $x$ and $y$ components. Similarly, a 3D
tensor can be drawn in 3-dimensional plane with $x$, $y$, and $z$
coordinates. Consider the following example.

\begin{figure}[H]
\centering
\tdplotsetmaincoords{70}{110}
\begin{tikzpicture}[
        tdplot_main_coords,
        axis/.style={->,thick},
        scale=0.5
    ]
    \foreach \x in {0,1,...,7}
        \foreach \y in {0,1,...,7} {
			\draw[gray,very thin]
                (\x,0,0) -- (\x,7.5,0);
			\draw[gray,very thin] (0,\y,0) -- (7.5,\y,0);
		}
    \foreach \y in {0,1,...,7}
        \foreach \z in {0,1,...,7} {
			\draw[gray,very thin]
                (0,\y,0) -- (0,\y,7.5);
			\draw[gray,very thin] (0,0,\z) -- (0,7.5,\z);
		}
    \foreach \z in {0,1,...,7}
        \foreach \x in {0,1,...,7} {
			\draw[gray,very thin]
                (0,0,\z) -- (7.5,0,\z);
			\draw[gray,very thin] (\x,0,0) -- (\x,0,7.5);
		}

	\draw[axis] (-2,0,0) -- (8,0,0) node[left] {$x$};
	\draw[axis] (0,-1,0) -- (0,8,0) node[right] {$y$};
	\draw[axis] (0,0,-1) -- (0,0,8) node[left] {$z$};

    \draw[->,very thick,Red,dashed] (0,0,0) -- (4,0,0);
    \draw[->,very thick,Blue,dashed] (4,0,0) -- (4,3,0);
    \draw[->,very thick,Green,dashed] (4,3,0) -- (4,3,5);
    \draw[->,very thick,black] (0,0,0) -- (4,3,5)
        node[right] {\small $\vec{v} = \langle 4,3,5 \rangle$};
\end{tikzpicture}
\end{figure}

Continuing, we can see different operations that vectors can perform.

\begin{definition}{}{}
\textbf{Vector addition} is a process of adding two vectors with
same dimension elementwise. \textbf{Scalar multiplication} is a
process of multiplying a scalar to individual elements.
\end{definition}

Let's take a look at quick examples. Consider the following equations.
\begin{align*}
    \langle 1, 2, 3 \rangle + \langle 4, 5, 6 \rangle
    &= \langle 1 + 4, 2 + 5, 3 + 9 \rangle
    = \langle 5, 7, 9 \rangle \\
    2 \langle 1, 2, 3 \rangle
    &= \langle 2 \cdot 1, 2 \cdot 2, 2 \cdot 3 \rangle
    = \langle 2, 4, 6 \rangle
\end{align*}
As you can see above, the operations are defined elementwise.
For vectors of different dimension, we say their sum is undefined.
Unlike addition and subtraction, multiplication is not defined
elementwise. When considering ``product'' we consider dot product
and matrix multiplication.

\begin{definition}{}{}
A \textbf{dot product} of two vectors is the sum of products of the
corresponding elements.
\end{definition}

Dot product will be quite important when we discuss attention and
transformers because it can tell us how ``related'' two vectors are.
Indeed, consider the vectors $A$ and $B$ with the angle between the
two vectors as $\theta$. The following equation holds.
\[
A \cdot B = ||A|| ||B|| \cos\theta
\]
Here, $||V||$ for $V \in \mathbb{R}^n$ represents the magnitude of the
vector, i.e. $\sqrt{\sum_{k=1}^n V_k^2}$. Let's take a look at a quick
example.
\begin{align*}
    \langle 1, 0 \rangle \cdot \langle 1, \sqrt{3} \rangle
    &= 1 \cdot 1 + 0 \cdot \sqrt{3} = 1 \\
    &= 1 \cdot 2 \cdot \cos60^\circ = 1
\end{align*}
One thing to keep in mind is that we have a special term for the
sign $||V||$.

\begin{definition}{}{}
The \textbf{vector norm} is the non-negative value that represents
the magnitude of a vector.
\end{definition}

With that in mind, let's discuss about matrices. There are so many
different and interesting applications of matrices such as solving
linear systems, but let's focus on the topics that will be essential
in machine learning.

There are different ways of representing a matrix $A$, but a common
way of writing is $[a_{ij}]$ where $a_{ij}$ are the entries in
$i^\text{th}$ row and $j^\text{th}$ column. A matrix $A_{m \times n}$
is said to have an \textbf{order} $m \times n$ if there are $m$ rows
and $n$ columns. For example, below is a $3 \times 2$ matrix.
\[
A =
\begin{bmatrix}
    a_{11} & a_{12} & a_{13} \\
    a_{21} & a_{22} & a_{23}
\end{bmatrix}
\]
Similar to vector addition, matrix addition and scalar multiplications
are defined elementwise. Moreover, the addition with matrices of
different order is undefined. Consider the following equations.
\begin{align*}
    [a_{ij}]_{m \times n} + [b_{ij}]_{m \times n}
    &= [a_{ij} + b_{ij}]_{m \times n} \\
    \alpha [a_{ij}]
    &= [\alpha a_{ij}]
\end{align*}
Now matrix multiplication is very different from the operations
above.

For matrices $[a_{ij}]_{m \times n}$ and $[b_{ij}]_{n \times p}$,
the product matrix $[c_{ij}]_{m \times p}$ is defined as the
following.
\[
c_{ij}
= \sum_{k=1}^{n} a_{ik} b_{kj}
\]
Below is an example.
\begin{align*}
\begin{bmatrix}
    1 & 2 & 3 \\
    4 & 5 & 6
\end{bmatrix}
\begin{bmatrix}
    1 & 2 \\
    3 & 4 \\
    5 & 6
\end{bmatrix}
&=
\begin{bmatrix}
    1 \cdot 1 + 2 \cdot 3 + 3 \cdot 5 &
    1 \cdot 2 + 2 \cdot 4 + 3 \cdot 6 \\
    4 \cdot 1 + 5 \cdot 3 + 6 \cdot 5 &
    4 \cdot 2 + 5 \cdot 4 + 6 \cdot 6 \\
\end{bmatrix} \\
&=
\begin{bmatrix}
    1 + 6 + 15  & 2 + 8 + 18 \\
    4 + 15 + 30 & 8 + 20 + 36
\end{bmatrix}
=
\begin{bmatrix}
    22 & 28 \\
    49 & 64
\end{bmatrix}
\end{align*}
Notice that the ``inner numbers'' of the orders must be the same
for the multiplication to be defined. Moreover, if we multiply
two matrices of orders $m \times n$ and $n \times p$ respectively,
the multiplication is defined since $n = n$. The product matrix will
now have the order $m \times p$. One more thing to note is that
we don't use the symbol $\cdot$ because that is reserved for dot
products. We can interpret matrix multiplication in many ways, but
let's now move on to a more abstract and advanced topics in linear
algebra for now.

\subsection{CONTINUE}






\end{document}


